\documentclass[10pt,a4paper]{amsart}
\usepackage[margin=19mm]{geometry}
\usepackage{amsmath,amssymb,mathtools}
\usepackage[scr=rsfs,cal=euler]{mathalfa}
\usepackage{xfrac}
\usepackage[unicode,hidelinks,bookmarks=false]{hyperref}
\newcommand{\ie}{i.e.\ }
\newcommand{\cf}{cf.\ }
\newcommand{\resp}{respectively\ }
\newcommand{\Subsection}{\subsection}
\providecommand{\nobreakdash}{\nobreak\hbox{-}}
\setlength{\emergencystretch}{2em}
\multlinegap0pt
\usepackage{bm}
\usepackage{bbm}
\usepackage{dsfont}
\usepackage{xspace}
\usepackage{cancel}
\usepackage{mathtools}
\usepackage{amsthm}
\makeatletter
\@ifundefined{lemma}{\newtheorem{lemma}{Lemma}}{}
\makeatother
\title[The Significance of Proposition II in Galois' M\'emoire, The]{The Significance of Proposition II in Galois' M\'emoire, The Origin of Galois Automorphisms}
\author{Math Dicker}
\date{Source: arXiv:2607.20147v1 (CC BY 4.0)}
\begin{document}
\maketitle

\noindent\emph{Source and licence.} The Significance of Proposition II in Galois' M\'emoire, The Origin of Galois Automorphisms. Version arXiv:2607.20147v1 (CC BY 4.0), distributed under the Creative Commons Attribution 4.0 International licence, as recorded in the arXiv licence metadata for this exact version and shown on the abstract page https://arxiv.org/abs/2607.20147v1. Full rights notice: licenses/arxiv-2607.20147-CC-BY-4.0.txt.\par\smallskip

\noindent\textit{Editorial scope.} Counted unit: the Lemma whose source label is \texttt{lem:een} in the article (source0.tex lines 183--186, source label \texttt{lem:een}), delivered together with the article's own complete original proof of it (source0.tex lines 188--323, one \texttt{proof} environment of 136 lines). No mathematics is written, repaired or re-derived: statement and proof are the article's own text line for line. Required context retained uncounted: none. Every definition, display and label of those lines is retained unchanged. Omitted from this package and not redistributed: 1--182, 187--187, 324--1558. Nothing inside a retained excerpt is omitted or altered: every retained body is delivered exactly as the article's lines are written, so no omission receipt of any kind accompanies this package. Citations: the retained excerpts carry no citation command at all, so no bibliography entry of the article is transcribed and no reference is redistributed. Reference adaptation: none was needed and none was made; every reference inside the delivered unit resolves inside it, so no unresolved reference or citation is produced. Printed slips kept literal: no printed word, symbol, hypothesis or number of the counted statement or of its proof was altered. Layout and class: the article's own class is \texttt{article} with packages including titling, babel, inputenc, graphicx, amsmath, amssymb, amsfonts, amsthm, geometry, float, enumitem, caption; the wrapper is the lane's pinned \texttt{amsart} editorial wrapper at 10pt on A4, which restates the theorem-like environments the retained text needs and adds the article's own macro definitions that the retained text needs (none), with extra wrapper packages (bm, bbm, dsfont, xspace, cancel, mathtools). The delivered numbering is the unit's own. Assets: no retained span contains a figure environment, a graphics-inclusion command, a PDF-page inclusion, a table or a TikZ picture; no image file is copied into this package, the package declares no asset dependency and no image file is redistributed with it. Licence: version arXiv:2607.20147v1 of this article is distributed under the Creative Commons Attribution 4.0 International licence, as recorded in the arXiv licence metadata for this exact version and shown on the abstract page https://arxiv.org/abs/2607.20147v1. A full rights notice is delivered at licenses/arxiv-2607.20147-CC-BY-4.0.txt. Characters: every retained excerpt is pure ASCII, so the delivered engine, XeLaTeX with its default Latin Modern text font, reports no missing character.
\begin{lemma}
	\label{lem:een}
	A set of permutations \(P=\{p_1,p_2,\dots,p_k\}\) is a \textbf{\textit{groupe de permutations}} if and only if \(P\) is a \textbf{right coset of a group}, if and only if \(P\) is a \textbf{left coset of a group}. In the former case, if \(P\) is a right coset of a group \(H\), then \(H\) is the corresponding \textbf{\textit{groupe de substitutions}}. In the latter case, if \(P\) is a left coset \(aH\), then \(aHa^{-1}\) is the corresponding \textbf{\textit{groupe de substitutions}}.
\end{lemma}

\begin{proof}[\normalfont\bfseries Proof]
	
	The last two equivalences are immediate. Indeed, if \(H\) is a group and \(a\) is a permutation, then
	\[
	aH=(aHa^{-1})a
	\]
	and
	\[
	Ha=a(a^{-1}Ha).
	\]
	
	First suppose that \(P\) contains the identity permutation \(1\). We then show that
	
	\[
	\{\text{\(P\) is a \textit{groupe de permutations} containing the identity}\}
	\Longleftrightarrow
	\{\text{\(P\) is a group}\}.
	\]
	
	Indeed, let \(x,y\in P\). Since \(P\) is a \textit{groupe de permutations}, there exists \(z\in P\) such that
	
	\[
	\left(
	\begin{array}{c}
		1\\
		x
	\end{array}
	\right)
	=
	\left(
	\begin{array}{c}
		y\\
		z
	\end{array}
	\right)
	=
	\left(
	\begin{array}{c}
		1\\
		zy^{-1}
	\end{array}
	\right).
	\]
	
	Hence \(x=zy^{-1}\), and therefore \(xy=z\in P\). Thus \(P\) is closed under composition.
	
	Conversely,
	
	\[
	\left(
	\begin{array}{c}
		p_i\\
		p_j
	\end{array}
	\right)
	=
	\left(
	\begin{array}{c}
		1\\
		p_jp_i^{-1}
	\end{array}
	\right),
	\]
	
	and since \(P\) is a group, \(p_jp_i^{-1}\in P\). Consequently, the defining condition of a \textit{groupe de permutations} is satisfied.
	
	Now suppose that \(P\) does not contain the identity permutation. Consider the set \(Pp_1^{-1}\). This is again a \textit{groupe de permutations}, but now containing the identity. By the previous argument it is therefore a group, say \(H\). Consequently,
	
	\[
	P=Hp_1,
	\]
	
	so that \(P\) is a right coset of \(H\).
	
	Conversely, suppose that \(P=Hp_1\) is a right coset of the group
	\(H=\{h_1,h_2,\dots,h_r\}\). Then
	
	\[
	\left(
	\begin{array}{c}
		h_ip_1\\
		h_jp_1
	\end{array}
	\right)
	=
	\left(
	\begin{array}{c}
		h_i\\
		h_j
	\end{array}
	\right),
	\]
	
	showing that \(Hp_1\) is a \textit{groupe de permutations}.
	
	Finally, suppose that \(P=Ha\) is a right coset of \(H\). Then
	
	\[
	\left(
	\begin{array}{c}
		h_ia\\
		h_ja
	\end{array}
	\right)
	=
	\left(
	\begin{array}{c}
		h_i\\
		h_j
	\end{array}
	\right),
	\]
	
	so that the associated substitutions are precisely those arising from \(H\).
	
	If, on the other hand, \(P=aH\) is a left coset of \(H\), then
	
	\[
	\left(
	\begin{array}{c}
		ah_i\\
		ah_j
	\end{array}
	\right)
	=
	\left(
	\begin{array}{c}
		ah_ia^{-1}\\
		ah_ja^{-1}
	\end{array}
	\right),
	\]
	
	and hence the associated substitutions are exactly those arising from the conjugate group \(aHa^{-1}\).
	
\end{proof}

\end{document}
